% !TeX root = ../sections/03_solutions.tex

\begin{solution} この問題では，\(x\) の作用が \[ xv=Av \] で定義されている。したがって， \[ x^2v=x(xv)=A(Av)=A^2v \] である。また，定数多項式 \(c\in\mathbb{R}\subset\mathbb{R}[x]\) は通常のスカラー倍として作用する。 よって，一般に \[ p(x)=c_0+c_1x+\cdots+c_nx^n\in\mathbb{R}[x] \] に対して， \[ p(x)v = c_0v+c_1Av+\cdots+c_nA^nv = p(A)v \] と計算すればよい。 \begin{enumerate} \item まず，\(\text{\rm (i)}\) を計算する。 \[ A \begin{pmatrix} 4\\ 5 \end{pmatrix} = \begin{pmatrix} 3 & 1\\ 0 & 2 \end{pmatrix} \begin{pmatrix} 4\\ 5 \end{pmatrix} = \begin{pmatrix} 3\cdot 4+1\cdot 5\\ 0\cdot 4+2\cdot 5 \end{pmatrix} = \begin{pmatrix} 17\\ 10 \end{pmatrix}. \] したがって， \[ x \begin{pmatrix} 4\\ 5 \end{pmatrix} = \begin{pmatrix} 17\\ 10 \end{pmatrix} \] である。ゆえに， \[ (2x+3) \begin{pmatrix} 4\\ 5 \end{pmatrix} = 2x \begin{pmatrix} 4\\ 5 \end{pmatrix} + 3 \begin{pmatrix} 4\\ 5 \end{pmatrix}. \] ここで， \[ 2x \begin{pmatrix} 4\\ 5 \end{pmatrix} = 2 \begin{pmatrix} 17\\ 10 \end{pmatrix} = \begin{pmatrix} 34\\ 20 \end{pmatrix} \] であり， \[ 3 \begin{pmatrix} 4\\ 5 \end{pmatrix} = \begin{pmatrix} 12\\ 15 \end{pmatrix} \] である。したがって， \[ (2x+3) \begin{pmatrix} 4\\ 5 \end{pmatrix} = \begin{pmatrix} 34\\ 20 \end{pmatrix} + \begin{pmatrix} 12\\ 15 \end{pmatrix} = \begin{pmatrix} 46\\ 35 \end{pmatrix}. \] よって， \[ \boxed{ (2x+3) \begin{pmatrix} 4\\ 5 \end{pmatrix} = \begin{pmatrix} 46\\ 35 \end{pmatrix} } \] である。 次に，\(\text{\rm (ii)}\) を計算する。 \[ A^2 = \begin{pmatrix} 3 & 1\\ 0 & 2 \end{pmatrix} \begin{pmatrix} 3 & 1\\ 0 & 2 \end{pmatrix} = \begin{pmatrix} 9 & 5\\ 0 & 4 \end{pmatrix}. \] したがって， \[ x^2 \begin{pmatrix} 3\\ -2 \end{pmatrix} = A^2 \begin{pmatrix} 3\\ -2 \end{pmatrix} = \begin{pmatrix} 9 & 5\\ 0 & 4 \end{pmatrix} \begin{pmatrix} 3\\ -2 \end{pmatrix}. \] よって， \[ A^2 \begin{pmatrix} 3\\ -2 \end{pmatrix} = \begin{pmatrix} 9\cdot 3+5\cdot(-2)\\ 0\cdot 3+4\cdot(-2) \end{pmatrix} = \begin{pmatrix} 17\\ -8 \end{pmatrix}. \] また， \[ 2 \begin{pmatrix} 3\\ -2 \end{pmatrix} = \begin{pmatrix} 6\\ -4 \end{pmatrix}. \] したがって， \[ (x^2-2) \begin{pmatrix} 3\\ -2 \end{pmatrix} = x^2 \begin{pmatrix} 3\\ -2 \end{pmatrix} - 2 \begin{pmatrix} 3\\ -2 \end{pmatrix}. \] ゆえに， \[ (x^2-2) \begin{pmatrix} 3\\ -2 \end{pmatrix} = \begin{pmatrix} 17\\ -8 \end{pmatrix} - \begin{pmatrix} 6\\ -4 \end{pmatrix} = \begin{pmatrix} 11\\ -4 \end{pmatrix}. \] よって， \[ \boxed{ (x^2-2) \begin{pmatrix} 3\\ -2 \end{pmatrix} = \begin{pmatrix} 11\\ -4 \end{pmatrix} } \] である。 \item \[ W= \left\{ \begin{pmatrix} a\\ 0 \end{pmatrix} \ \middle|\ a\in\mathbb{R} \right\} \] が \(V\) の \(\mathbb{R}[x]\)-部分加群であることを示す。 まず， \[ \begin{pmatrix} 0\\ 0 \end{pmatrix} \in W \] である。また， \[ \begin{pmatrix} a\\ 0 \end{pmatrix}, \begin{pmatrix} a'\\ 0 \end{pmatrix} \in W \] とすると， \[ \begin{pmatrix} a\\ 0 \end{pmatrix} + \begin{pmatrix} a'\\ 0 \end{pmatrix} = \begin{pmatrix} a+a'\\ 0 \end{pmatrix} \in W. \] さらに， \[ - \begin{pmatrix} a\\ 0 \end{pmatrix} = \begin{pmatrix} -a\\ 0 \end{pmatrix} \in W. \] したがって，\(W\) は \(V\) の加法部分群である。 次に，\(\mathbb{R}[x]\) の作用で閉じていることを示す。 任意に \[ \begin{pmatrix} a\\ 0 \end{pmatrix} \in W \] を取る。このとき， \[ x \begin{pmatrix} a\\ 0 \end{pmatrix} = A \begin{pmatrix} a\\ 0 \end{pmatrix} = \begin{pmatrix} 3 & 1\\ 0 & 2 \end{pmatrix} \begin{pmatrix} a\\ 0 \end{pmatrix} = \begin{pmatrix} 3a\\ 0 \end{pmatrix} \in W. \] したがって，\(x\) の作用で \(W\) は \(W\) の中に移る。 さらに，任意の \(k\in\mathbb{Z}_{\geq 0}\) に対して \[ x^k \begin{pmatrix} a\\ 0 \end{pmatrix} = A^k \begin{pmatrix} a\\ 0 \end{pmatrix} = \begin{pmatrix} 3^k a\\ 0 \end{pmatrix} \] が成り立つ。 実際，\(k=0\) のとき， \[ A^0 \begin{pmatrix} a\\ 0 \end{pmatrix} = I \begin{pmatrix} a\\ 0 \end{pmatrix} = \begin{pmatrix} a\\ 0 \end{pmatrix} = \begin{pmatrix} 3^0a\\ 0 \end{pmatrix}. \] また， \[ A^k \begin{pmatrix} a\\ 0 \end{pmatrix} = \begin{pmatrix} 3^ka\\ 0 \end{pmatrix} \] と仮定すると， \[ A^{k+1} \begin{pmatrix} a\\ 0 \end{pmatrix} = A \begin{pmatrix} 3^ka\\ 0 \end{pmatrix} = \begin{pmatrix} 3^{k+1}a\\ 0 \end{pmatrix}. \] よって，数学的帰納法により主張が従う。 ここで，任意の多項式 \[ p(x)=c_0+c_1x+\cdots+c_nx^n\in\mathbb{R}[x] \] を取る。このとき， \[ p(x) \begin{pmatrix} a\\ 0 \end{pmatrix} = \sum_{k=0}^{n}c_kx^k \begin{pmatrix} a\\ 0 \end{pmatrix}. \] 先ほどの計算より， \[ x^k \begin{pmatrix} a\\ 0 \end{pmatrix} = \begin{pmatrix} 3^ka\\ 0 \end{pmatrix} \] なので， \[ p(x) \begin{pmatrix} a\\ 0 \end{pmatrix} = \sum_{k=0}^{n} c_k \begin{pmatrix} 3^ka\\ 0 \end{pmatrix} = \begin{pmatrix} \sum_{k=0}^{n}c_k3^ka\\ 0 \end{pmatrix}. \] これは \(W\) の元である。したがって，任意の \(p(x)\in\mathbb{R}[x]\) と任意の \(w\in W\) に対して， \[ p(x)w\in W \] が成り立つ。 ゆえに， \[ \boxed{ W\text{ は }V\text{ の }\mathbb{R}[x]\text{-部分加群である。} } \] \item \(W\) が既約 \(\mathbb{R}[x]\)-加群であることを示す。 既約であるとは，\(W\) の \(\mathbb{R}[x]\)-部分加群が \[ 0 \quad\text{と}\quad W \] しかないことをいう。 \(U\) を \(W\) の \(\mathbb{R}[x]\)-部分加群とする。 すなわち， \[ U\subset W \] であり，\(U\) は \(\mathbb{R}[x]\) の作用で閉じているとする。 もし \[ U=0 \] ならば，これは自明な部分加群である。 そこで， \[ U\neq 0 \] とする。このとき，ある \[ 0\neq u\in U \] が存在する。 \(U\subset W\) であるから，ある \(a\in\mathbb{R}\) が存在して， \[ u= \begin{pmatrix} a\\ 0 \end{pmatrix} \] と書ける。さらに，\(u\neq 0\) であるから， \[ a\neq 0 \] である。 ここで，\(\frac{1}{a}\in\mathbb{R}\subset\mathbb{R}[x]\) を定数多項式とみなす。 \(U\) は \(\mathbb{R}[x]\)-部分加群なので， \[ \frac{1}{a}u\in U \] である。よって， \[ \frac{1}{a} \begin{pmatrix} a\\ 0 \end{pmatrix} = \begin{pmatrix} 1\\ 0 \end{pmatrix} \in U. \] 次に，任意の \(c\in\mathbb{R}\) を取る。 \(c\in\mathbb{R}\subset\mathbb{R}[x]\) を定数多項式とみなすと，\(U\) は \(\mathbb{R}[x]\)-部分加群だから， \[ c \begin{pmatrix} 1\\ 0 \end{pmatrix} \in U. \] すなわち， \[ \begin{pmatrix} c\\ 0 \end{pmatrix} \in U. \] これは任意の \(c\in\mathbb{R}\) について成り立つので， \[ W\subset U \] である。 一方で，初めから \[ U\subset W \] であった。したがって， \[ U=W \] である。 以上より，\(W\) の \(\mathbb{R}[x]\)-部分加群は \[ 0 \quad\text{または}\quad W \] に限られる。 ゆえに， \[ \boxed{ W\text{ は既約 }\mathbb{R}[x]\text{-加群である。} } \] \item 剰余加群 \(V/W\) における計算を行う。 剰余加群では， \[ v-v'\in W \] であるとき， \[ [v]=[v'] \] である。 任意の \[ \begin{pmatrix} a\\ b \end{pmatrix} \in V \] について， \[ \begin{pmatrix} a\\ b \end{pmatrix} - \begin{pmatrix} 0\\ b \end{pmatrix} = \begin{pmatrix} a\\ 0 \end{pmatrix} \in W \] である。したがって， \[ \left[ \begin{pmatrix} a\\ b \end{pmatrix} \right] = \left[ \begin{pmatrix} 0\\ b \end{pmatrix} \right]. \] つまり，\(V/W\) では第1成分は消え，第2成分だけが残る。 まず，\(x\) の作用を調べる。 \[ x \left[ \begin{pmatrix} 0\\ b \end{pmatrix} \right] = \left[ x \begin{pmatrix} 0\\ b \end{pmatrix} \right] = \left[ A \begin{pmatrix} 0\\ b \end{pmatrix} \right]. \] ここで， \[ A \begin{pmatrix} 0\\ b \end{pmatrix} = \begin{pmatrix} 3 & 1\\ 0 & 2 \end{pmatrix} \begin{pmatrix} 0\\ b \end{pmatrix} = \begin{pmatrix} b\\ 2b \end{pmatrix}. \] また， \[ \begin{pmatrix} b\\ 2b \end{pmatrix} - \begin{pmatrix} 0\\ 2b \end{pmatrix} = \begin{pmatrix} b\\ 0 \end{pmatrix} \in W \] である。したがって， \[ \left[ \begin{pmatrix} b\\ 2b \end{pmatrix} \right] = \left[ \begin{pmatrix} 0\\ 2b \end{pmatrix} \right]. \] ゆえに， \[ x \left[ \begin{pmatrix} 0\\ b \end{pmatrix} \right] = \left[ \begin{pmatrix} 0\\ 2b \end{pmatrix} \right]. \] したがって，剰余加群 \(V/W\) 上では，\(x\) は第2成分を \(2\) 倍する作用である。 まず，\(\text{\rm (i)}\) を計算する。 \[ (3x-2) \left[ \begin{pmatrix} 0\\ 3 \end{pmatrix} \right] = 3x \left[ \begin{pmatrix} 0\\ 3 \end{pmatrix} \right] - 2 \left[ \begin{pmatrix} 0\\ 3 \end{pmatrix} \right]. \] 先ほどの計算より， \[ x \left[ \begin{pmatrix} 0\\ 3 \end{pmatrix} \right] = \left[ \begin{pmatrix} 0\\ 6 \end{pmatrix} \right]. \] したがって， \[ 3x \left[ \begin{pmatrix} 0\\ 3 \end{pmatrix} \right] = 3 \left[ \begin{pmatrix} 0\\ 6 \end{pmatrix} \right] = \left[ \begin{pmatrix} 0\\ 18 \end{pmatrix} \right]. \] また， \[ 2 \left[ \begin{pmatrix} 0\\ 3 \end{pmatrix} \right] = \left[ \begin{pmatrix} 0\\ 6 \end{pmatrix} \right]. \] よって， \[ (3x-2) \left[ \begin{pmatrix} 0\\ 3 \end{pmatrix} \right] = \left[ \begin{pmatrix} 0\\ 18 \end{pmatrix} \right] - \left[ \begin{pmatrix} 0\\ 6 \end{pmatrix} \right]. \] したがって， \[ (3x-2) \left[ \begin{pmatrix} 0\\ 3 \end{pmatrix} \right] = \left[ \begin{pmatrix} 0\\ 12 \end{pmatrix} \right]. \] よって， \[ \boxed{ (3x-2) \left[ \begin{pmatrix} 0\\ 3 \end{pmatrix} \right] = \left[ \begin{pmatrix} 0\\ 12 \end{pmatrix} \right] } \] である。 次に，\(\text{\rm (ii)}\) を計算する。 \[ (2x^2+3) \left[ \begin{pmatrix} 0\\ 5 \end{pmatrix} \right] = 2x^2 \left[ \begin{pmatrix} 0\\ 5 \end{pmatrix} \right] + 3 \left[ \begin{pmatrix} 0\\ 5 \end{pmatrix} \right]. \] まず， \[ x \left[ \begin{pmatrix} 0\\ 5 \end{pmatrix} \right] = \left[ \begin{pmatrix} 0\\ 10 \end{pmatrix} \right]. \] さらに， \[ x^2 \left[ \begin{pmatrix} 0\\ 5 \end{pmatrix} \right] = x \left[ \begin{pmatrix} 0\\ 10 \end{pmatrix} \right] = \left[ \begin{pmatrix} 0\\ 20 \end{pmatrix} \right]. \] したがって， \[ 2x^2 \left[ \begin{pmatrix} 0\\ 5 \end{pmatrix} \right] = 2 \left[ \begin{pmatrix} 0\\ 20 \end{pmatrix} \right] = \left[ \begin{pmatrix} 0\\ 40 \end{pmatrix} \right]. \] また， \[ 3 \left[ \begin{pmatrix} 0\\ 5 \end{pmatrix} \right] = \left[ \begin{pmatrix} 0\\ 15 \end{pmatrix} \right]. \] よって， \[ (2x^2+3) \left[ \begin{pmatrix} 0\\ 5 \end{pmatrix} \right] = \left[ \begin{pmatrix} 0\\ 40 \end{pmatrix} \right] + \left[ \begin{pmatrix} 0\\ 15 \end{pmatrix} \right]. \] したがって， \[ (2x^2+3) \left[ \begin{pmatrix} 0\\ 5 \end{pmatrix} \right] = \left[ \begin{pmatrix} 0\\ 55 \end{pmatrix} \right]. \] よって， \[ \boxed{ (2x^2+3) \left[ \begin{pmatrix} 0\\ 5 \end{pmatrix} \right] = \left[ \begin{pmatrix} 0\\ 55 \end{pmatrix} \right] } \] である。 \end{enumerate} \end{solution}